\documentclass{article}

\usepackage{fancyhdr}
\usepackage{amsmath}
\allowdisplaybreaks
\usepackage{amsthm}
\usepackage{amsfonts}
\usepackage{tikz}
\usepackage[plain]{algorithm}
\usepackage{algpseudocode}
\usepackage{enumitem}
\usepackage{tikz}
\usepackage{pgfplots}
\pgfplotsset{compat=1.18}
\usepgfplotslibrary{groupplots}
\usepackage{comment}

\usetikzlibrary{automata,positioning}

%
% Basic Document Settings
%

\topmargin=-0.45in
\evensidemargin=0in
\oddsidemargin=0in
\textwidth=6.5in
\textheight=9.0in
\headsep=0.25in

\linespread{1.1}

\pagestyle{fancy}
\lhead{\hmwkAuthorName}
\chead{} % empty
\rhead{\hmwkHeaderTitle}
\cfoot{\thepage}

\renewcommand\headrulewidth{0.4pt}
\renewcommand\footrulewidth{0pt}

\setlength\parindent{0pt}

%
% Homework Details
%

\newcommand{\hmwkTitle}{Homework 7}
\newcommand{\hmwkDueDate}{Monday, February 23, 2026}
\newcommand{\hmwkClass}{ICMA 216 Calculus IIIA}
\newcommand{\hmwkClassTime}{Section 1}
\newcommand{\hmwkClassInstructor}{Ruth J. Skulkhu}
\newcommand{\hmwkAuthorName}{\textbf{Jiraroj Wiruchpongsanon (6781617)}}
\newcommand{\hmwkHeaderTitle}{ICMA 216 Calculus IIIA: Homework 7}

%
% Title Page
%

\title{
    \vspace{2in}
    \textmd{\textbf{\hmwkClass:\ \hmwkTitle}}\\
    \normalsize\vspace{0.1in}\small{Due\ on\ \hmwkDueDate}\\
    \vspace{0.1in}\large{\textit{\hmwkClassInstructor\ \hmwkClassTime}}
    \vspace{3in}
}

\author{\hmwkAuthorName}
\date{}

\renewcommand{\part}[1]{\textbf{\large Part \Alph{partCounter}}\stepcounter{partCounter}\\}

%
% Helper Commands
%

\newcounter{partCounter}
\newcommand{\solution}{\textbf{\large Solution}}

\newtheorem*{theorem}{Theorem}

\theoremstyle{remark}
\newtheorem*{remark}{Remark}

\begin{document}

\maketitle
\newpage

\section*{Problem 1 (13.3 Partial Derivatives)}
Let $f(x, y) = xe^{-y} + 5y$.
\begin{enumerate}[label=(\alph*)]
    \item Compute $f_x$ and $f_y$.
    \item Find the slope of the surface $z = f(x, y)$ in the $y$-direction at the point $(2, 5)$.
\end{enumerate}

\medskip
\solution

\begin{enumerate}[label=(\alph*)]
    \item for $f_x$, treat $y$ as constant, so
    \[
        f_x = \frac{\partial}{\partial x}(xe^{-y} + 5y) = e^{-y}
    \]

    and for $f_y$, treat $x$ as a constant
    \[
        f_y = \frac{\partial}{\partial y}(xe^{-y} + 5y) = -xe^{-y} + 5
    \]

    therefore,
    \[
        \boxed{f_x = e^{-y}, \qquad f_y = -xe^{-y} + 5}
    \]

    \item slope in $y$-direction, we do $f_y(2,5)$
    \[
        f_y(2,5) = -(2)e^{-5} + 5
    \]

    so the slope of the surface in the $y$-direction at $(2,5)$ is
    \[
        \boxed{-2e^{-5} + 5}
    \]
\end{enumerate}

\vspace{1em}
\hrulefill

\section*{Problem 2 (13.3 Partial Derivatives)}
Find all the second partial derivatives of the function
\begin{enumerate}[label=(\alph*)]
    \item $f(x, y) = x^4 - 3x^2y^3$.
    \item $f(x, y) = \ln(3x + 5y)$.
\end{enumerate}

\medskip
\solution

\begin{enumerate}[label=(\alph*)]
    \item
    \[
        f_x = 4x^3 - 6xy^3, \qquad f_y = -9x^2y^2
    \]

    therefore
    \[
        f_{xx} = 12x^2 - 6y^3
    \]
    \[
        f_{xy} = -18xy^2
    \]
    \[
        f_{yx} = -18xy^2
    \]
    \[
        f_{yy} = -18x^2y
    \]

    so,
    \[
        \boxed{
        \begin{aligned}
            f_{xx} &= 12x^2 - 6y^3, \\
            f_{xy} &= -18xy^2, \\
            f_{yx} &= -18xy^2, \\
            f_{yy} &= -18x^2y
        \end{aligned}}
    \]

    \item
    \[
        f_x = \frac{3}{3x+5y}, \qquad f_y = \frac{5}{3x+5y}
    \]

    now for
    \[
        f_{xx} = \frac{\partial}{\partial x}\left(\frac{3}{3x+5y}\right)
    \]

    let $u = 3x+5y \Rightarrow \frac{\partial u}{\partial x} = 3$

    \[
        3 \frac{\partial}{\partial x}(3x+5y)^{-1}
        = 3 \frac{\partial}{\partial u}(u^{-1}) \cdot \frac{\partial u}{\partial x}
    \]
    \[
        = 3(-1)u^{-2}\cdot 3 = \frac{-9}{(3x+5y)^2}
    \]

    and you do similar process with the other three, therefore
    \[
        \boxed{
        \begin{aligned}
            f_{xx} &= -\frac{9}{(3x+5y)^2}, \\
            f_{xy} &= -\frac{15}{(3x+5y)^2}, \\
            f_{yx} &= -\frac{15}{(3x+5y)^2}, \\
            f_{yy} &= -\frac{25}{(3x+5y)^2}
        \end{aligned}}
    \]
\end{enumerate}

\vspace{1em}
\hrulefill

\section*{Problem 3 (13.3 Partial Derivatives)}
If $z = f(x, y)$ with the relation $xy^2z^3 + x^3y^2z + 1 = x + y + z$. Find $\dfrac{\partial z}{\partial x}$ at $(1, 1, 1)$.

\medskip
\solution

we differentiate both side with respect to $x$:
\[
    \frac{\partial}{\partial x}(xy^2z^3) + \frac{\partial}{\partial x}(x^3y^2z) + \frac{\partial}{\partial x}(1) = \frac{\partial}{\partial x}(x+y+z)
\]

do it term-by-term, first one
\[
    \frac{\partial}{\partial x}(xy^2z^3) = y^2 \frac{\partial}{\partial x}(xz^3)
    = y^2\left(x(3z^2)\frac{\partial z}{\partial x} + z^3 \frac{\partial x}{\partial x}\right)
\]
\[
    = y^2\left(3xz^2\frac{\partial z}{\partial x} + z^3\right)
\]

and
\[
    \frac{\partial}{\partial x}(x^3y^2z) = y^2 \frac{\partial}{\partial x}(x^3z)
    = y^2\left(x^3 \frac{\partial z}{\partial x} + z \frac{\partial}{\partial x}(x^3)\right)
\]
\[
    = y^2\left(x^3\frac{\partial z}{\partial x} + 3x^2z\right)
\]

lastly,
\[
    \frac{\partial}{\partial x}(x+y+z) = 1 + 0 + \frac{\partial z}{\partial x}
\]

therefore
\[
    y^2\left(3xz^2\frac{\partial z}{\partial x} + z^3\right)
    +
    y^2\left(x^3\frac{\partial z}{\partial x} + 3x^2z\right)
    = 1 + \frac{\partial z}{\partial x}
\]

at $(x,y,z) = (1,1,1)$:
\[
    1\left(3(1)(1)^2\frac{\partial z}{\partial x} + 1\right)
    +
    1\left((1)^3\frac{\partial z}{\partial x} + 3(1)^2(1)\right)
    = 1 + \frac{\partial z}{\partial x}
\]

\[
    3\frac{\partial z}{\partial x} + 1 + \frac{\partial z}{\partial x} + 3
    = 1 + \frac{\partial z}{\partial x}
\]

\[
    4\frac{\partial z}{\partial x} + 4 = 1 + \frac{\partial z}{\partial x}
\]

\[
    3\frac{\partial z}{\partial x} = -3
\]

\[
    \boxed{\frac{\partial z}{\partial x} = -1}
\]

\vspace{1em}
\hrulefill

\section*{Problem 4 (13.4 Differentiability, Differentials, and Local Linearity)}
Use the linearization of $f(x, y) = \sqrt{x + 2y}$ to approximate $f(3.01, 2.96)$.

\medskip
\solution

we pick the base point at $(3,3)$, so

\[
    L(x,y) = f(3,3) + f_x(3,3)(x-3) + f_y(3,3)(y-3)
\]

solving each term:
\[
    f(3,3) = \sqrt{3 + 2(3)} = 3
\]

\[
    f_x = \frac{1}{2\sqrt{x+2y}}
    \quad \Rightarrow \quad
    f_x(3,3) = \frac{1}{2\sqrt{3 + 2(3)}} = \frac{1}{6}
\]

\[
    f_y = \frac{2}{2\sqrt{x+2y}} = \frac{1}{\sqrt{x+2y}}
    \quad \Rightarrow \quad
    f_y(3,3) = \frac{1}{\sqrt{3 + 2(3)}} = \frac{1}{3}
\]

therefore
\[
    L(x,y) = 3 + \frac{1}{6}(x-3) + \frac{1}{3}(y-3)
\]

\[
    f(3.01,2.96) \approx L(3.01,2.96)
\]

\[
    \approx 3 + \frac{1}{6}(3.01-3) + \frac{1}{3}(2.96-3)
\]

\[
    \boxed{\approx 2.9883}
\]

\vspace{1em}
\hrulefill

\section*{Problem 5 (13.4 Differentiability, Differentials, and Local Linearity)}
Find the differential of $z = y \cos(xy)$.

\medskip
\solution

as
\[
    dz = \frac{\partial z}{\partial x}dx + \frac{\partial z}{\partial y}dy
\]

\[
    \frac{\partial z}{\partial x}
    = y(-\sin(xy))y
    = -y^2\sin(xy)
\]

\[
    \frac{\partial z}{\partial y}
    = \cos(xy) + y[-\sin(xy)\cdot x]
    = \cos(xy) - xy\sin(xy)
\]

therefore
\[
    \boxed{dz = -y^2\sin(xy)\,dx + \left(\cos(xy) - xy\sin(xy)\right)\,dy}
\]

\vspace{1em}
\hrulefill

\section*{Problem 6 (13.4 Differentiability, Differentials, and Local Linearity)}
If $z = 5x^2 + y^2$ and $(x, y)$ changes from $(1, 2)$ to $(1.05, 2.1)$, find the value of $dz$.

\medskip
\solution

find $dz$ first:
\[
    dz = \frac{\partial z}{\partial x}dx + \frac{\partial z}{\partial y}dy
\]

\[
    dz = 10x\,dx + 2y\,dy
\]

\[
    dx = 1.05 - 1 = 0.05, \qquad dy = 2.1 - 2 = 0.1
\]

at $(x,y) = (1,2)$
\[
    dz = 10(1)(1.05 - 1) + 2(2)(2.1 - 2)
\]

\[
    = 10(0.05) + 4(0.1)
\]

\[
    \boxed{dz = 0.9}
\]

\end{document}
